% ============================================================================
% CHAPTER 2: Bewertungskatalog
% ============================================================================
\renewcommand{\chaptertag}{Bildfilter}
\chapter{Bewertungskatalog}

\vspace{-5pt}
{\color{maincolor}\rule{\textwidth}{0.4pt}}
\vspace{4pt}

{\color{maincolor}\textbf{\"Uberblick:}}
\begin{itemize}[leftmargin=*]
  \item Du kannst insgesamt \textbf{120 Punkte} erreichen, verteilt
    auf vier Bereiche.
  \item Den gr\"ossten Teil holst du dir mit den \textbf{Pflichtfiltern}.
    Daneben z\"ahlen dein \textbf{eigener Filter}, deine
    \textbf{Codequalit\"at} und eine ordentliche \textbf{Abgabe}.
  \item Die folgende Grafik zeigt, wieviele Punkte in jedem Bereich zu
    holen sind.
\end{itemize}

\vspace{8pt}

\begin{center}
\begin{tikzpicture}
  \def\W{15}     % Gesamtbreite in cm (entspricht textwidth)
  \def\hgt{1.4}  % Hoehe in cm
  % Anteile aus 120 Punkten: A=50 (0.4167), B=20 (0.1667), C=15 (0.125), D=35 (0.2917)
  \fill[maincolor]    (0,0)              rectangle (\W*0.4167, \hgt);
  \fill[maincolor!72] (\W*0.4167, 0)     rectangle (\W*0.5833, \hgt);
  \fill[maincolor!50] (\W*0.5833, 0)     rectangle (\W*0.7083, \hgt);
  \fill[maincolor!32] (\W*0.7083, 0)     rectangle (\W,        \hgt);
  % Beschriftungen
  \node[white, font=\bfseries]            at (\W*0.2083,\hgt*0.65) {A. Pflichtfilter};
  \node[white, font=\footnotesize]        at (\W*0.2083,\hgt*0.32) {50 Punkte};
  \node[white, font=\bfseries\small]      at (\W*0.5000,\hgt*0.65) {B. Eigen};
  \node[white, font=\scriptsize]          at (\W*0.5000,\hgt*0.32) {20 Pkt.};
  \node[white, font=\bfseries\scriptsize] at (\W*0.6458,\hgt*0.65) {C. Code};
  \node[white, font=\tiny]                at (\W*0.6458,\hgt*0.32) {15 Pkt.};
  \node[white, font=\bfseries]            at (\W*0.8542,\hgt*0.65) {D. Abgabe};
  \node[white, font=\footnotesize]        at (\W*0.8542,\hgt*0.32) {35 Punkte};
\end{tikzpicture}
\end{center}

\vspace{4pt}
{\color{maincolor}\rule{\textwidth}{0.4pt}}
\vspace{8pt}

\renewcommand{\arraystretch}{1.25}

% ============================================================================
\section{A. Pflichtfilter (50 Punkte)}

F\"ur jeden dieser Filter gibt es Punkte, wenn er korrekt funktioniert und
ein sichtbares Ergebnis liefert. Jeder Filter ist eine eigene Funktion.

\begin{center}
\begin{tabularx}{\textwidth}{>{\bfseries}p{4cm} X r}
\toprule
\tableheadercolor
Filter & Was er tut & Pkt. \\
\midrule
Kein Rot
  & Setzt den Rot-Wert aller Pixel auf 0.
  & 5 \\
\midrule
Invertieren
  & Aus jedem Wert wird $255 - $ Wert.
  & 7 \\
\midrule
Graustufen
  & Pro Pixel: \texttt{grau = (R + G + B) // 3}, neuer Pixel
    $(\text{grau}, \text{grau}, \text{grau})$.
  & 7 \\
\midrule
Helligkeit
  & Addiert eine Zahl $h$ auf jeden Farbkanal, klemmt die Werte
    auf $[0, 255]$.
  & 8 \\
\midrule
Schwellwert
  & Pixel \"uber einem Schwellwert $t$ werden weiss, alle anderen
    schwarz.
  & 7 \\
\midrule
Box-Blur
  & Pro Pixel: Durchschnitt aus dem Pixel selbst und seinen 8 Nachbarn
    ($3\times 3$-Fenster).
  & 10 \\
\midrule
Spiegeln
  & Spiegelt das Bild horizontal (linke und rechte Seite tauschen).
  & 6 \\
\bottomrule
\end{tabularx}
\end{center}

% ============================================================================
\section{B. Eigener Filter (20 Punkte)}

Erfinde und implementiere einen \textit{eigenen} Filter, oder w\"ahle
einen aus den Vorschl\"agen am Ende des Lernpfads.

\begin{center}
\begin{tabularx}{\textwidth}{>{\bfseries}p{4cm} X r}
\toprule
\tableheadercolor
Kriterium & Anforderung & Pkt. \\
\midrule
Funktion arbeitet korrekt
  & Der Filter produziert sichtbar das, was er soll.
  & 10 \\
\midrule
Eigenst\"andige Idee
  & Du kannst in zwei S\"atzen erkl\"aren, was dein Filter tut und warum.
  & 5 \\
\midrule
Beispielbild
  & Du zeigst ein Vorher/Nachher-Bild, auf dem der Effekt klar zu sehen ist.
  & 5 \\
\bottomrule
\end{tabularx}
\end{center}

% ============================================================================
\section{C. Codequalit\"at (15 Punkte)}

\begin{center}
\begin{tabularx}{\textwidth}{>{\bfseries}p{4cm} X r}
\toprule
\tableheadercolor
Kriterium & Anforderung & Pkt. \\
\midrule
Aussagekr\"aftige Namen
  & Variablen heissen wie das, was sie sind:
    \texttt{zeile}, \texttt{spalte}, \texttt{neuer\_pixel}, nicht
    \texttt{a}, \texttt{b}, \texttt{tmp}.
  & 5 \\
\midrule
Eine Funktion pro Filter
  & Jeder Filter ist eine eigene Funktion, nicht alles in einem grossen
    Block.
  & 5 \\
\midrule
Clamping korrekt
  & Bei \texttt{helligkeit} bleiben die Werte zwischen 0 und 255.
  & 3 \\
\midrule
Rand bewusst behandelt
  & Beim Box-Blur weisst du, was du am Bildrand machst (z.B.\ einfach
    weglassen).
  & 2 \\
\bottomrule
\end{tabularx}
\end{center}

% ============================================================================
\section{D. Abgabe (35 Punkte)}

\begin{center}
\begin{tabularx}{\textwidth}{>{\bfseries}p{4cm} X r}
\toprule
\tableheadercolor
Kriterium & Anforderung & Pkt. \\
\midrule
Code abgegeben
  & Alle Python-Dateien sind l\"auffahig und an die richtige Stelle hochgeladen.
  & 5 \\
\midrule
Drei Beispielbild-Paare
  & Drei Original/Filter-Paare als Bilddateien.
  & 5 \\
\midrule
Bericht
  & 2 bis 3 A4-Seiten: Was hast du gemacht? Was war schwierig?
    Was hat dich \"uberrascht?
  & 25 \\
\bottomrule
\end{tabularx}
\end{center}

\closechapter
